\begin{abstract}

Modelos de Linguagem de Grande Escala (\textit{Large Language Models}, LLMs) tornaram-se componentes centrais de aplicações modernas, mas ainda esbarram em uma limitação estrutural, a janela de contexto. O mecanismo de atenção apresenta complexidade quadrática no comprimento da sequência e, mesmo quando a janela é nominalmente ampla, o modelo tende a não utilizar toda a informação de forma eficaz, com fenômenos como \textit{lost in the middle}, \textit{context rot} e colapso de atenção. Este trabalho investiga estratégias não invasivas de extensão de contexto, aplicáveis a modelos acessados via API, sem retreinamento nem modificação arquitetural, e propõe um \textit{framework} que formula a compressão de \textit{prompt} como um problema de alocação de orçamento. A partir de uma revisão sistemática das técnicas existentes e de uma avaliação empírica de cinco estratégias em \textit{benchmarks} de contexto longo, o trabalho formaliza a seleção de compressores como um problema da mochila de múltipla escolha. Dado um \textit{prompt} particionado, um orçamento efetivo de contexto e uma família de compressores parametrizados, o modelo escolhe, para cada partição, a opção de compressão que maximiza a informação preservada sob a restrição de orçamento. Uma triagem adicional compara LLMLingua-GPT2, LLMLingua-2, Selective Context e uma aproximação CPC-MiniLM em cinco modelos e dois orçamentos de contexto. CPC-MiniLM obteve a maior média entre os quatro compressores em oito das dez combinações, enquanto LLMLingua-2 liderou nas duas combinações do Qwen 3 30B-A3B. No Gemma 4 26B com 32768 \textit{tokens}, CPC-MiniLM alcançou 0.573 e apresentou a maior média entre as nove técnicas comparadas. Esses resultados sustentam a seleção inicial de CPC-MiniLM, LLMLingua-2 e Selective Context como opções do MCKP. A avaliação completa do \textit{framework} proposto é apresentada como plano experimental reprodutível.

\end{abstract}

\begin{foreignabstract}

Large Language Models (LLMs) have become central components of modern applications, yet they still face a structural limitation, the context window. The attention mechanism has quadratic complexity in the sequence length and, even when the window is nominally large, the model tends not to use all of the available information effectively, exhibiting phenomena such as lost in the middle, context rot, and attention collapse. This work investigates non-invasive strategies for context extension, applicable to models accessed through an API, without retraining or architectural modification, and proposes a framework that formulates prompt compression as a budget allocation problem. Building on a systematic review of existing techniques and an empirical evaluation of five strategies on long-context benchmarks, the work formalizes compressor selection as a multiple-choice knapsack problem. Given a partitioned prompt, an effective context budget, and a family of parameterized compressors, the model selects, for each partition, the compression option that maximizes the preserved information under the budget constraint. An additional screening compares LLMLingua-GPT2, LLMLingua-2, Selective Context, and a CPC-MiniLM approximation across five models and two context budgets. CPC-MiniLM achieved the highest average among the four compressors in eight of the ten model-budget combinations, whereas LLMLingua-2 led in both Qwen 3 30B-A3B combinations. On Gemma 4 26B with 32,768 tokens, CPC-MiniLM reached 0.573 and achieved the highest average among all nine techniques. These results support the initial selection of CPC-MiniLM, LLMLingua-2, and Selective Context as MCKP options. The complete evaluation of the proposed framework is presented as a reproducible experimental plan.

\end{foreignabstract}
